\section{Domain Instantiations under One Runtime Contract}
\label{sec:ieee-domains}

The key test of a universal safety layer is not whether every domain shares the same
dynamics or controller. It is whether the same runtime contract survives domain change
without losing semantic meaning. ORIUS evaluates that question across Battery, AV, and
Healthcare under explicit evidence boundaries.

\begin{table*}[t]
\centering
\caption{ORIUS domain instantiations under a shared runtime contract.}
\label{tab:ieee-domain-rows}
\small
\begin{tabularx}{\textwidth}{@{}p{0.14\textwidth}p{0.19\textwidth}p{0.16\textwidth}p{0.15\textwidth}X@{}}
\toprule
\textbf{Domain} & \textbf{Degraded-observation hazard} & \textbf{Safety object} & \textbf{Current tier} & \textbf{Exact closure boundary}\\
\midrule
Battery energy storage & Dispatch legal on observed state but unsafe on true state of charge or reserve surface & State-of-charge and dispatch envelope preservation & Primary validation domain & Deepest closure surface in the current stack.\\
Autonomous vehicles & Action proposed on stale or optimistic headway / collision state & TTC and entry-barrier preservation under repair & Bounded nuPlan replay domain & Bounded to completed all-zip grouped nuPlan replay artifacts under the current TTC plus predictive-entry-barrier contract.\\
Healthcare monitoring & Clinical alerting or intervention released on degraded vital telemetry & Vital-alert envelope and bounded intervention discipline & Bounded MIMIC monitoring domain & Bounded to retrospective monitoring and alert-release semantics; no live clinical deployment claim.\\
\bottomrule
\end{tabularx}
\end{table*}

\subsection{Autonomous vehicles}
The AV row is defended under a bounded safety object: time-to-collision and predictive
entry-barrier preservation when the perception or state-estimation surface degrades.

\paragraph{AV results ($n = 1{,}531{,}104$ nuPlan steps).}
The nominal controller (no runtime layer) produces a baseline TSVR of
$\mathbf{28.93\%}$ --- more than one-quarter of AV steps involve an action that
is legal on the perceived state but unsafe on the true physical state.
With ORIUS, TSVR falls to $\mathbf{0.016\%}$, a $99.94\%$ relative reduction.
Certificate-valid release rate is $99.99\%$; runtime latency p95 is $0.74$\,ms.
Compared to the always-brake degenerate fail-safe ($69{,}226$ useful-work units),
ORIUS delivers $116{,}192$ units while keeping excess TSVR over that fail-safe
at $\approx 0$.  The fallback activation rate of $17.4\%$ confirms that the
T6/T7 temporal certificate and brake-hold contract are actively governing the
AV release surface.  The evidence is bounded to the TTC plus predictive-entry-barrier
contract and does not claim full autonomous-driving field closure.


\subsection{Healthcare monitoring}
Healthcare matters because degraded observation can suppress intervention rather than
only create excessive actuation.  ORIUS still fits this setting because the runtime layer
is defined over admissible decisions, fallback policies, and certificate semantics.

\paragraph{Healthcare results ($n = 137{,}262$ MIMIC steps).}
The nominal monitoring controller produces a baseline TSVR of $\mathbf{19.45\%}$;
with ORIUS this falls to $\mathbf{0.000\%}$, complete elimination of true-state
violations.  Certificate-valid release rate is $100.0\%$; latency p95 is $0.70$\,ms.
Compared to the always-alert degenerate fail-safe (zero useful monitoring work),
ORIUS retains $142{,}767$ units of clinically useful output while maintaining zero
excess TSVR, a non-vacuous utility-preserving safety result.  The healthcare evidence
is bounded to retrospective MIMIC monitoring semantics; it does not claim prospective
clinical deployment or clinical-decision-support approval.

\subsection{Future additional domains}
The active branch does not defend named additional domains beyond Battery, AV, and
Healthcare. Future extensions remain part of the architectural argument, but they are
not part of the current empirical claim surface until they satisfy the same replay,
runtime, governance, and bounded-prose obligations as the active three-domain program.
